\section{Background}\label{sec:background}

\subsection{Smart Home Simulation}

Simulating smart home environments requires modeling three
interdependent domains: the \emph{physical space} (rooms, furniture,
appliances), the \emph{occupant behavior} (movement, activities,
routines), and the \emph{device ecosystem} (sensors, actuators,
communication protocols, cloud services).  Existing tools address
these in isolation.

\textbf{3D Embodied Simulation.}
Habitat~\cite{habitat19iccv,habitat2,habitat3}, AI2-THOR~\cite{ai2thor},
and iGibson~\cite{igibson} provide photorealistic 3D environments with
physics engines, but focus on embodied AI navigation tasks rather
than IoT device modeling.  Habitat~3.0 introduced humanoid avatars
with realistic motion and multi-agent coordination, making it
suitable as a foundation for occupant simulation.

\textbf{IoT Testbeds and Simulators.}
IoTSim~\cite{iotsim}, FIESTA-IoT~\cite{fiesta-iot}, and SWoTSuite~\cite{swotsuite} model
device networks but lack 3D spatial grounding.  They represent devices
as abstract message endpoints, missing firmware-level behaviors such as
sensor drift, interrupt timing, and flash-storage wear.

\textbf{Activity Generation.}
Traditional approaches use Markov chains~\cite{markov-activity}, stochastic Petri
nets~\cite{petri-net-activity}, or agent-based models~\cite{abm-activity} to generate occupant
activities.  These methods require extensive manual parameterization
and struggle to produce the contextual variety observed in real
households with diverse occupant demographics.

\subsection{Firmware-in-the-Loop Testing}

Firmware-in-the-loop (FIL) testing runs actual device firmware on
emulated hardware to catch bugs that unit tests and hardware-in-the-loop
setups miss.  QEMU~\cite{qemu} provides full-system ARM emulation and has
been used for IoT firmware analysis~\cite{firmadyne,avatar}.  \sys{}
packages QEMU instances in Docker containers, each running a compiled
ARM binary with simulated GPIO, UART, and network peripherals.

\subsection{LLMs for Behavioral Modeling}

Recent work has explored LLMs as world models~\cite{llm-world-model} and behavior
planners~\cite{llm-planner}.  Park et al.~\cite{generative-agents} demonstrated
that LLM-driven agents can produce believable daily routines in a
virtual town.  \sys{} adapts this insight to the smart home domain:
instead of open-ended sandbox behavior, we use structured prompts to
generate \emph{activity schedules}---timestamped sequences of
activities with room assignments and device interactions---that can be
validated against real-world activity datasets.

\subsection{SmartThings Ecosystem}

Samsung SmartThings is a leading IoT cloud platform supporting
hundreds of device types via its Schema Connector and Device
Profiles.  The Schema Connector protocol defines a webhook-based
lifecycle (discovery, state refresh, command handling) through which
third-party devices integrate with the SmartThings app.  \sys{}'s
bridge implements this protocol, allowing simulated devices to appear
alongside real hardware in a user's SmartThings account.
